% arara: lualatex: {shell: yes, options: ["--output-directory=build"]}
% arara: biber: {options: ["-input-directory=build", "-output-directory=build"]}
% arara: lualatex: {shell: yes, options: ["--output-directory=build"]}

\documentclass[english, frametitlelogo, mainlogoleft, handout]{beamer}
%\documentclass[english, frametitlelogo, mainlogoleft, article]{beamer}
%\documentclass[english, frametitlelogo, mainlogoleft, notes]{beamer}

% -- Extra symbols
\usepackage{amssymb}
\usepackage{amsthm}
\usepackage{empheq}
\DeclarePairedDelimiter{\ceil}{\lceil}{\rceil}
\DeclarePairedDelimiter{\floor}{\lfloor}{\rfloor}
\DeclarePairedDelimiter{\braces}{\{}{\}}
\DeclareMathOperator{\lcm}{\text{\slshape\mdseries\sffamily lcm}}
% override some commands so they don't get rm fonts
\let\max\relax
\DeclareMathOperator*{\max}{\text{\slshape\mdseries\sffamily max}}
\let\min\relax
\DeclareMathOperator*{\min}{\text{\slshape\mdseries\sffamily\text{min}}}
\let\mod\relax
\DeclareMathOperator{\mod}{\text{\slshape\mdseries\sffamily\text{mod}}}
\let\gcd\relax
\DeclareMathOperator{\gcd}{\text{\slshape\mdseries\sffamily\text{gcd}}}
\let\exp\relax
\DeclareMathOperator{\exp}{\text{\slshape\mdseries\sffamily\text{exp}}}
% user defined commands
\newrobustcmd*\bigO{\mathcal{O}}
\newrobustcmd*\mathdash{\text{-}}

% -- Bibliography and citations
\usepackage[
	backend = biber,
	style = numeric,
	sorting = ynt,
    maxbibnames = 8,
    maxcitenames = 2,
	%alldates=iso
	]{biblatex}
%\usepackage{fvextra}
\usepackage[autostyle]{csquotes}

% -- Glossary and nomenclature (usage with bib2gls)
\usepackage[record, symbols]{glossaries-extra}

% --  Floats/images
\graphicspath{{graphics/}}
\usepackage{caption}
\usepackage{subcaption}
\usepackage{flafter}
\usepackage{tabularx}
\usepackage{multirow}
\usepackage{ragged2e}
\usepackage{multicol}

% -- Drawing schemes and plotting
\usepackage{tikz}
\usepackage{pgfplots}
\usetikzlibrary{shapes, shapes.misc, shapes.symbols, positioning, decorations.pathreplacing, external, arrows.meta, patterns, plotmarks, hobby}
\tikzset{>=Latex} % solid arrow tip shape
\pgfplotsset{compat=1.16}
\pgfplotsset{table/search path = {data}}
\usepgfplotslibrary{statistics}
\pgfkeys{/pgf/number format/.cd,1000 sep={}} % remove comma as thousands separator

% -- Pseudocode and source code
\usepackage{algorithm}
\usepackage{algpseudocode}
\usepackage{verbatim}

\renewcommand{\textproc}{\texttt} % typeset algorithmic function names
% enable linking to algorithm line numbers. credit to https://tex.stackexchange.com/questions/148977/incorrect-reference-to-a-line-in-algorithmic-using-hyperref
\newcommand\plabel[1]{\phantomsection\label{#1}}
% define the for each loop in pseudocode. credit to https://tex.stackexchange.com/questions/149162/how-can-i-define-a-foreach-loop-on-the-basis-of-the-existing-forall-loop
\algnewcommand\algorithmicforeach{\textbf{for each}}
\algdef{S}[FOR]{ForEach}[1]{\algorithmicforeach\ #1\ \algorithmicdo}

% -- Tables
\usepackage{booktabs}
\usepackage{makecell}

% -- Misc
\usepackage[disable]{todonotes}
\usepackage{subfiles}

% -- Extra math - units
\usepackage[binary-units=true]{siunitx}

% --- Specific to the DIMA-scheduler thesis document ---

% -- custom colors
\colorlet{colorA}{blue!50!white}
\colorlet{colorB}{green!50!white}
\colorlet{colorC}{magenta!80!white}
\colorlet{colorD}{orange!90}
\colorlet{swred}{red!75!black!70!white}
\colorlet{hwgray}{gray!10!white}

% -- styles
\tikzstyle module_style_1=[hwgray, dashed, fill=cyan, fill opacity=0.1]
\tikzstyle module_style_2=[hwgray, dashed, fill=magenta, fill opacity=0.1]
\tikzstyle{decision} = [diamond, draw, text badly centered, node distance=3cm, inner sep=0pt, text width=1.8cm]
\tikzstyle{block} = [rectangle, draw, text centered, rounded corners, minimum height=1cm, inner sep=0pt, text width=3cm] 
\tikzstyle{line} = [draw,->]
\tikzstyle{file} = [tape, tape bend top=none, tape bend height=2mm, draw, minimum height=1cm, text width=2.6cm, inner sep=0pt, text centered]

% -- scripts
\directlua{require("scripts/tableformat.lua")}
\directlua{require("scripts/loadlines.lua")}

%\usepackage{pgfpages}
%\setbeameroption{show notes}
%\setbeameroption{show notes on second screen=right}

\addbibresource{bibliography.bib}

\mode<presentation> { 
    \usetheme{Lisboa} 
    %\usecolortheme{dolphin}
}

% \\\vspace{0.6cm} \small\itshape \href{mailto:{\myemail}}{\myemail}

\newrobustcmd*\myemail{joaomfgoncalves@tecnico.ulisboa.pt}
\newrobustcmd*\mytitle{Partition Scheduling in Distributed \linebreak Integrated Modular Avionics}
\newrobustcmd*\mytitleshort{Partition Scheduling in DIMA}
\newrobustcmd*\myauthor{João Miguel Fonseca Gonçalves}
\newrobustcmd*\myauthorshort{João Gonçalves}
\newrobustcmd*\myinstitution{Instituto Superior Técnico}
\newrobustcmd*\myinstitutionshort{IST-UL}
\newrobustcmd*\mydate{November 29, 2019}
\newrobustcmd*\mydegree{MEAer}
\newrobustcmd*\mynumber{81040}

\title{\mytitleshort}
\date{\mydate}
\institute{\myinstitutionshort}
\author{\myauthorshort}

\hypersetup{
    pdfsubject = {Artificial Intelligence, Operations Research, Avionics},
    pdfkeywords = {Arinc 653, Integrated Modular Avionics, Scheduling, Optimization},
} 

\begin{document}

\begin{frame}

    \begin{tikzpicture}[node distance=1pt, shift={(0,0)}]

        \node at (0, 3) {};

        \fill[fill=white] (0,0) rectangle node[color=ist-cyan, text width=120mm, align=center] (A) {\LARGE\mytitle} ++(329pt, 2);

        \node[below=0.7cm of A] (B) {\Large\myauthor};
        \node[below=0.005cm of B, color=ist-cyan] (C) {\small\href{mailto:\myemail}{\myemail}};
        \node[below=0.2cm of C] (D) {\myinstitution};
        \node[below=0.2cm of D] (E) {\footnotesize\mydate};

		\pgfmathsetmacro\r{-4.5}
        \node at (2.5, \r) {\includegraphics[height=2.5cm]{IST_A_CMYK_POS.pdf}};
        \node at (8.5, \r) {\includegraphics[width=3.7cm]{ULISBOA_HORIZONTAL_CMYK.pdf}};
%        \node at (0,0) {\includegraphics[height=2.5cm]{IST_A_CMYK_POS.pdf}};
%        \node at (4,0) {\includegraphics[width=3.7cm]{ULISBOA_HORIZONTAL_CMYK.pdf}};
    \end{tikzpicture}

    \note<1->{
        \begin{itemize}
            \item Aviónica Modular Integrada e Distribuída
        \end{itemize}
    }
\end{frame}

\begin{frame}
    \frametitle{Outline}
    \begin{minipage}{0.49\linewidth}
        \tableofcontents[sections=1-2]
    \end{minipage}
    \begin{minipage}{0.49\linewidth}
        \tableofcontents[sections=3-5]
    \end{minipage}
\end{frame}

\section{Motivation}
\subsection{Avionic Architectures}

\begin{frame}
    \frametitle{Avionic Architectures}
    \begin{columns}[T]
        \begin{column}{0.4\linewidth}

            Federated Avionics

            {\small
            \begin{itemize}
                \item Dedicated hardware for each application
            \end{itemize}
            }

            \vspace{20pt}

            IMA

            {\small
            \vspace{3pt}
            \begin{itemize}
                \item Resource sharing between unrelated applications
                \item Active research area
            \end{itemize}
            }
        \end{column}
        \begin{column}{0.6\linewidth}
            \centering
            \resizebox{\linewidth}{!}{%\tikzstyle my grid=[style=help lines, color=blue!50, step=0.5, very thin, dashed]
\begin{tikzpicture}
	%\draw[style=my grid] (-1.5, -0.5) grid (10,9);

	
	\draw (7.5, 9) node {\LARGE Federated Avionics};
	\foreach \x in {0, 1, ..., 7} {
        \pgfmathsetmacro\r{\x*2}
		\pgfmathsetmacro\s{\r+1}
        \pgfmathsetmacro\l{\r+0.5}
        \draw[black!50!gray, thick, fill=hwgray] (\r,6) rectangle node[black] {HW} (\s,7);
        \draw[black!50!gray, thick, fill=swred] (\r,7) rectangle node[black] {SW} (\s, 8);
        \draw[<->] (\l, 6) -- (\l, 5.5);
    }
	\draw (16, 7) node {$\ldots$};
    \draw[dashed, rounded corners] (0, 5) rectangle node {} (15, 5.5);

	
	\draw (7.5, 3) node {\LARGE Integrated Modular Avionics};
	
	\foreach \p in {3, 8} {
        \foreach \x in {0, 1, 2, 3} {
            \pgfmathsetmacro\l{\x+\p}
            \pgfmathsetmacro\r{\x+\p+1}
            \draw[black!50!gray, thick, fill=swred] (\l,1) rectangle node[black] {SW} (\r,2);
        }
        \pgfmathsetmacro\r{\p+4}
        \pgfmathsetmacro\m{\p+2}
        \draw[black!50!gray, thick, fill=hwgray] (\p, 0) rectangle node[black] {HW} (\r, 1);
        \draw[<->] (\m, 0) -- (\m, -1);
	}
	\draw (13, 1) node {$\ldots$};
	
    \draw[dashed, rounded corners] (3, -1.5) rectangle node {} (12, -1);
	
\end{tikzpicture}
}
        \end{column}
    \end{columns}
    \note<1->{
        \begin{itemize}
            \item Aviónica Federada
            \item Aviónica Modular Integrada
            \item Época onde se deu a transição - década 90
            \item Robustez e segurança devida à separação física de aplicações
            \item Modularizar -- integrar aplicações não relacionadas em hardware partilhado
        \end{itemize}
    }
\end{frame}


\begin{frame}
    \frametitle{Distributed Integrated Modular Avionics}
    \begin{columns}[T]
        \begin{column}{0.4\linewidth}
            Benefits:
            \begin{itemize}
                \item<1-> Maintains safety
                \vspace{10pt}
                \item<2-> Efficient usage of resources
                \vspace{10pt}
                \item<3-> Reduced space, weight and power
                \vspace{10pt}
                \item<4-> Uses open standards
                \vspace{10pt}
                \item<5-> Enables the involvement of smaller players
            \end{itemize}
        \end{column}
        \begin{column}{0.6\linewidth}
            \centering
            \resizebox{\linewidth}{!}{%\tikzstyle my grid=[style=help lines, color=blue!50, step=0.5, very thin, dashed]
\begin{tikzpicture}

	\draw (7.5, 3) node {\large Distributed IMA};
	
	\foreach \p in {3, 8} {
        \foreach \x in {0, 1, 2, 3} {
            \pgfmathsetmacro\l{\x+\p}
            \pgfmathsetmacro\r{\x+\p+1}
            \draw[black!50!gray, thick, fill=swred] (\l,1) rectangle node[black] {SW} (\r,2);
        }
        \pgfmathsetmacro\r{\p+4}
        \pgfmathsetmacro\m{\p+2}
        \draw[black!50!gray, thick, fill=hwgray] (\p, 0) rectangle node[black] {HW} (\r, 1);
        \draw[<->] (\m, 0) -- (\m, -1);
	}
	
    \draw[dashed, rounded corners] (3, -1.5) rectangle node {\itshape Shared Avionic Bus} (12, -1);

    \foreach \p in {0, 4} {
        \pgfmathsetmacro\l{\p+5}
        \draw[black!50!gray, thick, fill=hwgray] (\l,-3.2) rectangle node[black] {RDC} ++(1.1,1);
    }
    \foreach \p in {5.55, 9.55} {
        \draw[<->] (\p, -2.2) --++(0, 0.7);
        \draw[<->] (\p, -4) --++(0, 0.8);
    }
    \foreach \p in {0, 4} {
        \pgfmathsetmacro\l{\p+4.5}
        \draw[black!50!gray, thick, rounded corners, fill=ist-cyan!50!gray] (\l,-4.8) rectangle node[black] {Payload} ++(2,0.8);
    }
	
\end{tikzpicture}
}
        \end{column}
    \end{columns}
    \note<1->{
        
        \begin{itemize}
            \item Aviónica Modular Integrada e Distribuída
            \item Remote Data Concentrators (Concentrador remoto de dados) -- dispositivo miniaturizado que faz o IO com a rede aviónica e os periféricos
            \item Sumário: Isolamento por software vs. por hardware
        \end{itemize}
    }
\end{frame}


\begin{frame}
    \frametitle{Strong Partitioning}
    How is this achievable?

    \vfill

    \textbf{Space Partitioning}
    \uncover<2->{
    \begin{itemize}
        \item Statically isolated memory for each application
    \end{itemize}
    }

    \vspace{40pt}

    \textbf{Time Partitioning}
    \uncover<3->{
    \begin{itemize}
        \item<3-> Guaranteed processor time for each application
        \item<4-> Deterministic scheduling 
    \end{itemize}
    }
    \note<1->{
        \begin{itemize}
            \item Particionamento robusto
            \item Uma aplicação não pode aceder a dados de outra, quer em memória ou em trânsito
        \end{itemize}
    }
\end{frame}


\begin{frame}
    \frametitle{Arinc 653}

    \begin{itemize}
        \item Avionics Application Software Standard Interface
    \end{itemize}

    \begin{center}
        \resizebox{0.7\linewidth}{!}{%\tikzstyle my grid=[style=help lines, color=blue!50, step=1, very thin, dashed]
\begin{tikzpicture}[scale=1.0]

	%\draw[style=my grid] (-1.5, -0.5) grid (21.5,10.5);

	\colorlet{grayshade}{gray!40!white}

	\draw[very thick] (0,0) rectangle (20,10.75);
	\draw (10, 11.5) node {\huge Core Processing Module};
	
	\draw[thick, fill=hwgray] (0.5, 0.5) rectangle (19.5, 1.5);
	\draw (10,1) node {\Large HW};
	
	\draw[<->, very thick] (10, 1.5) -- (10, 2.5);
	
	\draw[thick, fill=colorB] (0.5,2.5) rectangle (19.5, 5);
	\draw (10,3.75) node {\Large Arinc 653 compliant RTOS};
	
	\draw[thick, fill=colorA] (0.5, 5) rectangle (19.5, 5.75);
	\draw (10, 5.375) node {\Large APEX};
	
	\foreach \x in {0.5, 5.5, 10.5, 15.5} {
		\pgfmathsetmacro\r{\x+4}
		\pgfmathsetmacro\c{\x+2}
		\draw[thick, fill=swred] (\x, 6.5) rectangle (\r, 9.5);
		\draw[<->, very thick] (\c, 6.5) -- (\c, 5.75);
	}
	\draw (2.5, 10) node {\Large Partition 1};
	\draw (7.5, 10) node {\Large Partition 2}; 
	\draw (12.5, 10) node {\Large Partition 3}; 
	\draw (17.5, 10) node {\Large Partition 4}; 	
	
	\pgfmathsetmacro{\ybottom}{7.2}
	\pgfmathsetmacro{\ytop}{8.5}
	\tikzstyle process=[ellipse, fill=grayshade, draw=none, minimum height=0.2, minimum width=0.6]
	
	\foreach \x in {1.5, 6.5, 11.5, 16.5} {
		\draw (\x, \ytop) node[style=process]  {\large P1};
	}
	\foreach \x in {1.5, 6.5, 16.5} {
		\draw (\x, \ybottom) node[style=process]  {\large P2};
	}
	\foreach \x in {1.5,  16.5} {
		\pgfmathsetmacro\y{\x+2}
		\draw (\y, \ytop) node[style=process]  {\large P3};
	}
	\foreach \x in {1.5} {
		\pgfmathsetmacro\xx{\x+2}
		\draw (\xx, \ybottom) node[style=process]  {\large P4};
	}	
\end{tikzpicture}
}
    \end{center}
    \note<1->{
        \begin{itemize}
            \item Interface entre Aplicações e o Sistema Operativo de Tempo Real
            \item Partições é o nome que se dá às aplicações
        \end{itemize}
    }
\end{frame}


\subsection{Partition Schedules}
\begin{frame}
    \frametitle{Partition Schedules}
    \begin{columns}[T]
        \begin{column}{0.4\linewidth}
            Characteristics
            \begin{itemize}
                \item<1-> Static
                \item<2-> Strictly periodic
                \item<3-> Functional requirements:
                \uncover<3->{
                    \begin{itemize}
                        \item Inter-partition communication
                        \item Access to resources
                        \item Redundancy configuration
                    \end{itemize}
                }
            \end{itemize}

            \vfill
            \uncover<4->{
            Status
            \begin{itemize}
                \item<4-> Manual, iterative process
                \item<5-> \textit{In-house} solutions
            \end{itemize}
            }
        \end{column}
        \begin{column}{0.6\linewidth}
            \centering
            \resizebox{0.9\linewidth}{!}{\begin{tikzpicture}[scale=1.0]
	%\tikzstyle my grid=[style=help lines, color=blue!50, step=1, very thin, dashed]
	%\draw[style=my grid] (-1, -0.5) grid (25,3);	

	\foreach \x in {0, 6} {
		\pgfmathsetmacro\r{\x+1}
		\draw[gray, fill=colorA] (\x, 1) rectangle node[black] {P1} (\r, 1.5);
		\draw[gray, fill=colorA] (\x, -0.7) rectangle node[black] {P1} (\r, -0.2);
	}
	\foreach \x in {7} {
		\pgfmathsetmacro\r{\x+2}
		\draw[gray, fill=colorB] (\x, 0.5) rectangle node[black] {P2} (\r, 1);
		\draw[gray, fill=colorB] (\x, -0.7) rectangle node[black] {P2} (\r, -0.2);
	}
	\foreach \x in {1} {
		\pgfmathsetmacro\r{\x+3}
		\draw[gray, fill=colorD] (\x, 0) rectangle node[black] {P3} (\r, 0.5);
		\draw[gray, fill=colorD] (\x, -0.7) rectangle node[black] {P3} (\r, -0.2);
	}
	\foreach \x/\s in {4/2, 9/3} {
		\pgfmathsetmacro\r{\x+\s}
		\draw[gray, fill=hwgray] (\x, -0.7) rectangle node[black] {Idle} (\r, -0.2);
	}
	
	\foreach \x in {0, 12} {
		\draw[dashed, gray, thin] (\x, -1.4) -- (\x, 1.7);
	}

	%\draw[|->, very thick] (0,-1.1) -- (24, -1.1) node[below] {$t$};
	
	\draw[<->, thick] (0, 1.7) -- node[above] {\large Partition 1 period} (6, 1.7);
	\draw[<->, thick] (0, -1.4) -- node[below] {\Large Major Time Frame} (12, -1.4);

\end{tikzpicture}
}
        \end{column}
    \end{columns}
    \note<1->{
        \begin{itemize}
            \item Estático -- Definido em configuração, não altera durante a execução, e possivelmente durante o tempo de vida do sistema
            \item Estritamente periódico -- As partições repetem-se com um periodo, e os tempos de inicio são separados exatamente pelo periodo, sem flexibilidade
            \item Diferença com os problemas de escalonamento usuais
            \item Tempo -- largos meses a anos
        \end{itemize}
    }
\end{frame}

\begin{frame}
    \frametitle{Partition Schedules -- Examples}
    \centering
    \resizebox{0.9\linewidth}{!}{\begin{tikzpicture}[scale=0.9]
	%\tikzstyle my grid=[style=help lines, color=blue!50, step=1, very thin, dashed]
	%\draw[style=my grid] (-1, -0.5) grid (26,6);	

	\pgfmathsetmacro\h{0.5};
	
	\foreach \x in {1,...,23} {
		\draw[dashed, gray, thin] (\x, -0.6) -- (\x, 5.3);
	}
	
	\pgfmathsetmacro\y{4};
	\foreach \x in {0, 6, 12, 18} {
		\pgfmathsetmacro\r{\x+1}
		\pgfmathsetmacro\t{\y+\h};
		\draw[gray, fill=colorA] (\x, \y) rectangle node[black] {$p_1$} (\r, \t);
		\draw[gray, fill=colorA] (\x, 0) rectangle node[black] {$p_{1}$} (\r, \h);
	}
	\pgfmathsetmacro\y{3};
	\foreach \x in {7, 19} {
		\pgfmathsetmacro\r{\x+2}
		\pgfmathsetmacro\t{\y+\h};
		\draw[gray, fill=colorB] (\x, \y) rectangle node[black] {$p_2$} (\r, \t);
		\draw[gray, fill=colorB] (\x, 0) rectangle node[black] {$p_2$} (\r, \h);
	}
	\pgfmathsetmacro\y{2};
	\foreach \x in {1, 9, 17} {
		\pgfmathsetmacro\r{\x+1}
		\pgfmathsetmacro\t{\y+\h};
		\draw[gray, fill=colorC] (\x, \y) rectangle node[black] {$p_3$} (\r, \t);
		\draw[gray, fill=colorC] (\x, 0) rectangle node[black] {$p_3$} (\r, \h);
	}
	\pgfmathsetmacro\y{1};
	\foreach \x in {2, 14} {
		\pgfmathsetmacro\r{\x+3}
		\pgfmathsetmacro\t{\y+\h};
		\draw[gray, fill=colorD] (\x, \y) rectangle node[black] {$p_4$} (\r, \t);
		\draw[gray, fill=colorD] (\x, 0) rectangle node[black] {$p_4$} (\r, \h);
	}
	\foreach \x/\s in {5/1, 10/2, 13/1, 21/3} {
		\pgfmathsetmacro\r{\x+\s}
		\draw[gray, fill=hwgray] (\x, 0) rectangle node[black] {Idle} (\r, \h);
	}	

	
	\foreach \x in {0, 24} {
		\draw[thick, dashed] (\x, -0.6) -- (\x, 5.3);
	}
	
	\draw[<->, thick] (0,-0.6) -- node[below] {\Large MTF} (24, -0.6);
	

\end{tikzpicture}
}

    \vspace{1cm}
    \centering
    \resizebox{0.9\linewidth}{!}{\begin{tikzpicture}[scale=1.0]
	%\tikzstyle my grid=[style=help lines, color=blue!50, step=1, very thin, dashed]
	%\draw[style=my grid] (-1, -0.5) grid (21,3);	

	\foreach \x in {0, 4, 8, 12} {
		\pgfmathsetmacro\r{\x+1}
		\draw[gray, fill=colorA] (\x, 0.5) rectangle node[black] {P1} (\r, 1);
		\draw[gray, fill=colorA] (\x, -0.7) rectangle node[black] {P1} (\r, -0.2);
	}
	\foreach \x in {1, 9} {
		\pgfmathsetmacro\r{\x+2.75}
		\pgfmathsetmacro\l{\x+1.5}
		\draw[gray, fill=colorB] (\x, 0.5) -- (\x, 0) -- (\r, 0) arc[start angle=-90, end angle=90, radius=0.25] -- cycle;
		\draw[gray, fill=colorB] (\x, -0.2) -- (\x, -0.7) -- (\r, -0.7) arc[start angle=-90, end angle=90, radius=0.25] -- cycle;
		\draw (\l, 0.25) node {P2};
		\draw (\l, -0.45) node {P2};
	}

    \draw[gray, fill=colorB] (0, 1.5) rectangle node[black] {P2} (4.5, 2);
    \draw node at (4.7, 1.75) {?};

	\foreach \x in {5, 13} {
		\pgfmathsetmacro\s{\x+0.25}
		\pgfmathsetmacro\l{\x+0.5}
		\pgfmathsetmacro\r{\x+1}
		\draw[gray, fill=colorB] (\s, 0.5) arc[start angle=-270, end angle=-90, radius=0.25] -- (\r, 0) -- (\r, 0.5) -- cycle;
		\draw[gray, fill=colorB] (\s, -0.2) arc[start angle=-270, end angle=-90, radius=0.25] -- (\r, -0.7) -- (\r, -0.2) -- cycle;
		\draw (\l, 0.25) node {P2};
		\draw (\l, -0.45) node {P2};
	}
	\foreach \x/\s in {6/2, 14/2} {
		\pgfmathsetmacro\r{\x+\s}
		\draw[gray, fill=hwgray] (\x, -0.7) rectangle node[black] {Idle} (\r, -0.2);
	}
	
	\foreach \x in {0, 8, 16} {
		\draw[dashed, gray, thin] (\x, -1.4) -- (\x, 1.7);
	}
	
	\draw[<->, thick] (0, -1.4) -- node[below] {\large MTF} (8, -1.4);
	\draw[<->, thick] (8, -1.4) -- node[below] {\large MTF} (16, -1.4);
	%\draw[|->, very thick] (0,-1.1) -- (24, -1.1) node[below] {$t$};

\end{tikzpicture}
}
    \note<1->{
        \begin{itemize}
            \item Extenções ao problema
            \item Divisão de cada instância da partição em múltiplas janelas (tratado separadamente)
            \item Motivação na figura de baixo
        \end{itemize}
    }
\end{frame}

\subsection{Objectives}

\begin{frame}
    \frametitle{Objectives}

    Challenges:
    \vspace{5pt}
    \begin{itemize}
        \item<1-> NP-completeness
        \item<2-> Defining the system requirements
        \item<3-> Reconfiguration is lengthy and expensive
    \end{itemize}

    \vspace{20pt}

    Objectives:
    \vspace{5pt}
    \begin{itemize}
        \item<4-> Comprehensive mathematical description
        \item<5-> Automate schedule generation and validation
        \item<6-> Find flexible solutions
    \end{itemize}

\end{frame}

\section{Partition Scheduling Model}
\subsection{Variables}

\begin{frame}
    \frametitle{Partition Scheduling Model}
    \tableofcontents[currentsection,hideallsubsections]
\end{frame}

\begin{frame}
    \frametitle{The Partition Scheduling Problem}

    \begin{minipage}[t]{0.62\linewidth}
    Schedule $N_p$ partitions:

    \begin{itemize}
        \item[$T_i$] -- partition $i$ period
        \item[$e_i$] -- partition $i$ execution time (WCET)
        \item[$s_i$] -- partition $i$ memory requirement
    \end{itemize}

    \vspace{5pt}
    in $N_c$ modules:

    \begin{itemize}
        \item[$S_m$] -- module $m$ memory capacity
    \end{itemize}
    \end{minipage}
    \begin{minipage}[t]{0.37\linewidth}
        \vspace{5pt}
        \uncover<2->{
        Find, for each partition $i$:
        \begin{itemize}
            \item[$f_i$] -- assigned module
            \item[$t_i$] -- starting offset
        \end{itemize}
        }
    \end{minipage}

    \vfill
    \centering
    \resizebox{0.9\linewidth}{!}{\begin{tikzpicture}[scale=0.9]
	%\tikzstyle my grid=[style=help lines, color=blue!50, step=1, very thin, dashed]
	%\draw[style=my grid] (-1, -0.5) grid (26,6);	

	\pgfmathsetmacro\h{0.5};
	
	\foreach \x in {1,...,23} {
		\draw[dashed, gray, thin] (\x, -0.6) -- (\x, 5.3);
	}
	
	\pgfmathsetmacro\y{4};
	\foreach \x in {0, 6, 12, 18} {
		\pgfmathsetmacro\r{\x+1}
		\pgfmathsetmacro\t{\y+\h};
		\draw[gray, fill=colorA] (\x, \y) rectangle node[black] {$p_1$} (\r, \t);
		\draw[gray, fill=colorA] (\x, 0) rectangle node[black] {$p_{1}$} (\r, \h);
	}
	\pgfmathsetmacro\y{3};
	\foreach \x in {7, 19} {
		\pgfmathsetmacro\r{\x+2}
		\pgfmathsetmacro\t{\y+\h};
		\draw[gray, fill=colorB] (\x, \y) rectangle node[black] {$p_2$} (\r, \t);
		\draw[gray, fill=colorB] (\x, 0) rectangle node[black] {$p_2$} (\r, \h);
	}
	\pgfmathsetmacro\y{2};
	\foreach \x in {1, 9, 17} {
		\pgfmathsetmacro\r{\x+1}
		\pgfmathsetmacro\t{\y+\h};
		\draw[gray, fill=colorC] (\x, \y) rectangle node[black] {$p_3$} (\r, \t);
		\draw[gray, fill=colorC] (\x, 0) rectangle node[black] {$p_3$} (\r, \h);
	}
	\pgfmathsetmacro\y{1};
	\foreach \x in {2, 14} {
		\pgfmathsetmacro\r{\x+3}
		\pgfmathsetmacro\t{\y+\h};
		\draw[gray, fill=colorD] (\x, \y) rectangle node[black] {$p_4$} (\r, \t);
		\draw[gray, fill=colorD] (\x, 0) rectangle node[black] {$p_4$} (\r, \h);
	}
	\foreach \x/\s in {5/1, 10/2, 13/1, 21/3} {
		\pgfmathsetmacro\r{\x+\s}
		\draw[gray, fill=hwgray] (\x, 0) rectangle node[black] {Idle} (\r, \h);
	}	

	\draw[<->, thick] (0,-0.6) -- node[below] {$H_m$} (24, -0.6);
	
	\foreach \x in {0, 24} {
		\draw[thick, dashed] (\x, -0.6) -- (\x, 5.3);
	}
	
	\draw[thick, ->] (0, 1.25) -- node[above] {$t_4$} (2, 1.25);
	\draw[thick, ->] (0, 2.25) -- node[above] {$t_3$} (1, 2.25);
	\draw[thick, ->] (0, 3.25) -- node[above] {$t_2$} (7, 3.25);
	
	\draw[thick, <->] (2, 5.3) -- node[above] {$e_4$} (5, 5.3);
	\draw[thick, <->] (6, 5.3) -- node[above] {$e_1$} (7, 5.3);
	
	\draw[thick, <->] (9, 5.3) -- node[above] {$T_3$} (17, 5.3);
	

\end{tikzpicture}
}
    \note<1->{
        \begin{itemize}
            \item WCET -- Tempo de Execução no pior caso
            \item $t_i$ Desfazamento inicial
            \item $f_i$ -- módulo onde é escalonada a partição
            \item Cada módulo tem um destes escalonamentos
        \end{itemize}
    }
\end{frame}

\subsection{Distribution Constraints}

\begin{frame}
    \frametitle{Distribution Constraints}

    \begin{itemize}
        \item \textbf{Domain}  -- partitions can be scheduled in only some modules
        \item \textbf{Exclusion} -- some partitions must be scheduled in distinct modules
        \item \textbf{Inclusion} -- some partitions must be scheduled in the same module
        \item \textbf{Memory} -- a module's memory must not be exceeded
        \item \textbf{Uniqueness} -- a partition can only be scheduled in one module
    \end{itemize}
    \note<1->{
        \begin{itemize}
            \item Restições que afetam a disposição de partições entre os módulos
        \end{itemize}
    }
\end{frame}

\subsection{Communication Constraints}

\begin{frame}
    \frametitle{Inter-partition Communication}

    \begin{itemize}
        \item $E_{i,j}$ -- communication \textquote{chains}
        \item $τ_{m,n}$ -- inter-module communication delay
        \item Synchronous communication
    \end{itemize}

    \only<1>{\centering\resizebox{0.9\linewidth}{!}{\begin{tikzpicture}[scale=1.0]
	%\tikzstyle my grid=[style=help lines, color=blue!50, step=1.0, very thin, dashed]

	\draw[style=module_style_1] (0,2) rectangle (12, 4) node[below left, black, opacity=1.0] {$c_1$};
	\draw[style=module_style_2] (0,0) rectangle (12, 2) node[below left, black, opacity=1.0] {$c_2$};
	\foreach \x in {0,...,12} {
		\draw[color=gray, very thin, dashed] (\x, 0) --++(0, 5);
	}

    \draw[->, thick] (0,4) -- (12,4) node[above, yshift=0.1mm] {$t$};
    \draw[->, thick] (0,0) -- (12,0) node[below, yshift=-0.1mm] {$t$};

	
	\draw[gray, fill=colorA] (0, 2.75) rectangle node[black] {$p_1$} (2, 3.25);
	
	\draw[gray, fill=colorB] (4, 0.75) rectangle node[black] {$p_2$} (5, 1.25);
	\draw[gray, fill=colorB] (10, 0.75) rectangle node[black] {$p_2$} (11, 1.25);
	
	\draw[->] (2,3) -- (2,2) -- node[above] {$\tau_{1,2}$} (3,2);
	\draw[->] (3,2) -- (3,1) -- (4,1);
	
	\draw[<->] (0, 5) -- node[above] {$E_{1,2}$} ++(5,0);

\end{tikzpicture}
}}
    \only<2>{\centering\resizebox{0.9\linewidth}{!}{\begin{tikzpicture}[scale=1.0]
	%\tikzstyle my grid=[style=help lines, color=blue!50, step=1.0, very thin, dashed]

	\draw[style=module_style_1] (0,2) rectangle (12, 4) node[below left, black, opacity=1.0] {$c_1$};
	\draw[style=module_style_2] (0,0) rectangle (12, 2) node[below left, black, opacity=1.0] {$c_2$};
	\foreach \x in {0,...,12} {
		\draw[color=gray, very thin, dashed] (\x, 0) --++(0, 5);
	}

    \draw[->, thick] (0,4) -- (12,4) node[above, yshift=0.1mm] {$t$};
    \draw[->, thick] (0,0) -- (12,0) node[below, yshift=-0.1mm] {$t$};

	
	\draw[gray, fill=colorA] (0, 2.75) rectangle node[black] {$p_1$} (2, 3.25);
	
	\draw[gray, fill=colorB] (3, 0.75) rectangle node[black] {$p_2$} (4, 1.25);
	\draw[gray, fill=colorB] (9, 0.75) rectangle node[black] {$p_2$} (10, 1.25);
	
	\draw[->] (2,3) -- (2,2) -- node[above] {$\tau_{1,2}$} (3,2);
	\draw[->] (3,2) -- (3,1);
	
	\draw[<->] (0, 5) -- node[above] {$E_{1,2}$} ++(10,0);

\end{tikzpicture}
}}

    \note<1->{
        \begin{itemize}
            \item Comunicação modelada como cadeias
            \item Menor separação entre inicio e fim de duas partições é limitada
        \end{itemize}
    }

\end{frame}

\subsection{Timing Constraints}

\begin{frame}
    \frametitle{Timing Constraints}

    \begin{itemize}
        \item No temporal overlap between two partitions in the same module
    \end{itemize}

    \vspace{12pt}

    Theorem \cite{korst1992periodic}:

    \begin{align*}
        & g_{i,j} = {gcd}\braces*{T_i,T_j} \\[5pt]
        & e_i ≤ {mod}\braces*{t_j-t_i, g_{i,j}} ≤ T_j-e_j
    \end{align*}

    \vfill

    {\centering\resizebox{\linewidth}{!}{\begin{tikzpicture}[scale=1.3]
	%\tikzstyle my grid=[style=help lines, color=blue!50, step=1, very thin, dashed]
	%\draw[style=my grid] (0, -0.5) grid (12,1.5);
	
	\draw[->, thin, swred] (0, -1.0) -- (0, 1.3) node[above, black] {\scriptsize MTF};
    \draw[->, thin, swred] (12, -1.0) -- (12, 1.3) node[above, black] {\scriptsize MTF};
	
	\draw[thick, |->] (0,0) -- (12,0) node[right] {$t$};
	\foreach \x in {1, 4, 7} {
		\draw[gray, thin, dashed] (\x, -1.0) -- (\x, 1.3);
	}
	
	\foreach \x in {1, 7} {
		\pgfmathsetmacro\r{\x+2}
		\draw[gray, thin, fill=colorA] (\x, 0) rectangle node[black] {$p_i$} (\r,0.5);
	}
	\draw[gray, thin, fill=colorB] (4, 0) rectangle node[black] {$p_j$} (5, 0.5);
	
	\draw[|->] (0, -0.5) -- node[above] {$t_i$} (1, -0.5);
	\draw[|->] (0,-0.8) -- node[below] {$t_j$} (4,-0.8);
	
	\draw[<->] (1, 1.1) -- node[above] {$\text{mod}\braces{t_j-t_i, g_{i,j}}$} (4, 1.1);
	\draw[<->] (4, 0.8) -- node[above] {$\text{mod}\braces{t_i-t_j, g_{i,j}}$} (7, 0.8);

\end{tikzpicture}
}}

\end{frame}

\begin{frame}
    \frametitle{Timing Constraints}

    \begin{itemize}
        \item Multiple windows of execution require an execution penalty, $ε$
        \item Response time, $r_i$, is bounded
    \end{itemize}

    \vfill

    {\centering\resizebox{\linewidth}{!}{\begin{tikzpicture}[scale=0.85]
	%\tikzstyle my grid=[style=help lines, color=blue!50, step=1, very thin, dashed]
    %\draw[style=my grid] (-1.5, -2) grid (21,4);

    \draw[->, thick] (0,0) -- (21,0) node [below=2pt] {$t$};
	\draw[->, thin, swred] (0, -1.6) -- (0, 2.3) node[above, black] {\footnotesize MTF};
	\draw[->, thin, swred] (20, -1.6) -- (20, 2.3) node[above, black] {\footnotesize MTF};
	
	% \draw[gray, dashed] (10, -1.6) -- (10, 2.3);
	\draw[gray, dashed] (3, -0.4) -- (3, 2.3);
	\draw[gray, dashed] (13, -1.6) -- (13, 2.3);
	% \draw[gray, dashed] (17.5, -1.2) -- (17.5, 1.5);
    

	
	\draw[black!50!gray, thick, fill=colorA] (3,0) rectangle (8,1);
	\draw[black!50!gray, thick, fill=colorA] (13,0) -- (13,1) -- (15.5,1) arc[start angle=90, end angle=-90, radius=0.5] -- cycle; 
	%\draw[black!50!gray, thick, fill=colorA] (13,0) rectangle (16,1);
	\draw[black!50!gray, thick, fill=colorA] (18,0) arc[start angle=270, end angle=90, radius=0.5] -- (20,1) -- (20,0) -- cycle;
	%\draw[black!50!gray, thick, fill=colorA] (18,0) rectangle (20,1);
	
	\draw (5.5, 0.5) node {$\lambda_{i,1,1}$};
	\draw (14.5, 0.5) node {$\lambda_{i,2,1}$};
	\draw (19, 0.5) node {$\lambda_{i,2,2}$};

	\draw[<->] (3, 1.3) -- node[above] {$e_{i,1,1}$} (8, 1.3);
	\draw[<->] (13, 1.3) -- node[above] {$e_{i,2,1}$} (16, 1.3);
	\draw[<->] (17.5, 1.1) -- node[above] {$e_{i,2,2}$} (20, 1.1);
	\draw[<->] (17.5, 1.5) -- node [above] {$\varepsilon$} (18, 1.5);
	\draw[|->] (13, -1.6) -- (20, -1.6) node[right] {$r_{i}$};

\end{tikzpicture}
}}

    \note<1->{
        \begin{itemize}
            \item No caso do modelo com múltiplas janelas
            \item Tempo de resposta limitado
            \item Execução total deve dar o WCET, acrescido do tempo de context-switch
            \item Context Switch - troca de contexto, involve restaurar caches, toma tempo 
        \end{itemize}
    }
\end{frame}

\subsection{Optimization}

\begin{frame}
    \frametitle{Optimizing Potential Execution}

    \begin{itemize}
         \item Partition 1 cannot increase its execution
    \end{itemize}
    \vspace{10pt}
    {\centering\resizebox{\linewidth}{!}{\begin{tikzpicture}[scale=1.0]
	%\tikzstyle my grid=[style=help lines, color=blue!50, step=1, very thin, dashed]
	%\draw[style=my grid] (-1, -0.5) grid (17,2);	
	
	\draw[->, thick] (0,0) -- (16,0) node[below] {$t$};
	
	%\foreach \x in {5, 8} {
	%	\draw[gray, dashed, thin] (\x, 0) -- (\x, 1.5);
	%}
	
	\foreach \x in {0, 8} {
		\draw[gray, fill=colorA] (\x, 0) rectangle node[black] {P1} ++(2.5, 0.5);
	}
	\draw[gray, fill=colorB] (2.5,0) rectangle node[black] {P2} ++(1, 0.5);

	
\end{tikzpicture}
}}

    \vspace{42pt}

    \begin{itemize}
         \item It is possible to increase execution times without rescheduling
    \end{itemize}
    \vspace{10pt}
    {\centering\resizebox{\linewidth}{!}{\begin{tikzpicture}[scale=1.0]
	%\tikzstyle my grid=[style=help lines, color=blue!50, step=1, very thin, dashed]
	%\draw[style=my grid] (-1, -0.5) grid (17,2);	
	
	\draw[->, thick] (0,0) -- (16,0) node[below] {$t$};
	
	%\foreach \x in {5, 8} {
	%	\draw[gray, dashed, thin] (\x, 0) -- (\x, 1.5);
	%}
	
	\foreach \x in {0, 8} {
		\draw[gray, fill=colorA] (\x, 0) rectangle node[black] {P1} ++(2.5, 0.5);
	}
	\draw[gray, fill=colorB] (5, 0) rectangle node[black] {P2} ++(1, 0.5);

	
	
\end{tikzpicture}
}}
    \note<1->{
        \begin{itemize}
            \item Não só encontrar soluções viáveis, como mais favoráveis
            \item Objetivo é aumentar a flexibilidade do sistema
            \item Permitir acrescentar funcionalidade sem refazer o escalonamento
        \end{itemize}
    }
\end{frame}

\begin{frame}
    \frametitle{Optimizing Potential Execution}

    {\centering\resizebox{0.9\linewidth}{!}{\begin{tikzpicture}[scale=1.0]
	%\tikzstyle my grid=[style=help lines, color=blue!50, step=1, very thin, dashed]
	%\draw[style=my grid] (-1, -0.5) grid (17,2);	
	
	\draw[->, thick] (0,0) -- (16,0) node[below] {$t$};
	
	%\foreach \x in {5, 8} {
	%	\draw[gray, dashed, thin] (\x, 0) -- (\x, 1.5);
	%}
	
	\foreach \x in {0, 8} {
		\pgfmathsetmacro\r{\x+2.5}
		\draw[gray, fill=colorA] (\x, 0) rectangle node[black] {$p_1$} (\r, 0.5);
	}
	\draw[gray, fill=colorB] (5,0) rectangle node[black] {$p_2$} (6, 0.5);
	
	\foreach \x in {2.5, 10.5} {
		\pgfmathsetmacro\r{\x+2.5}
		\draw[dashed, gray, fill=colorA, fill opacity=0.3] (\x, 0) rectangle (\r, 0.5);
	}
	\draw[dashed, gray, fill=colorB, fill opacity=0.3] (6, 0) rectangle (7, 0.5);
	
	\draw[<->] (0,-0.2) -- node[below] {$e_1$} (2.5, -0.2);
	\draw[<->] (5,-0.2) -- node[below] {$e_2$} (6, -0.2);
	\draw[<->] (0, 0.8) -- node[above] {$α e_1$} (5, 0.8);
	\draw[<->] (5, 0.8) -- node[above] {$α e_2$} (7, 0.8);
	
	
\end{tikzpicture}
}}

    \vspace{20pt}

    Objective: maximize $α$

    \vspace{5pt}

    \begin{itemize}
        \item[$α$] -- \textquote{Minimum factor that scales all execution windows without overlaps}
    \end{itemize}
    \note<1->{
        \begin{itemize}
            \item Factor proporcional à duração da tarefa
            \item Horário é viável com $α≥1$
        \end{itemize}
    }
\end{frame}

\section{Methodology}

\begin{frame}
    \frametitle{Partition Scheduling Model}
    \tableofcontents[currentsection,hideallsubsections]
\end{frame}

\subsection{Mixed Integer Linear Programming}

\begin{frame}{Mixed Integer Linear Program}

    \begin{itemize}
        \item<1-> Optimization problem with both integer and real-valued variables
        \item<2-> Branch and bound / branch and cut algorithms
        \item<3-> Yields an optimal solution
    \end{itemize}

    \vfill

    {\centering\resizebox{0.95\linewidth}{!}{\begin{minipage}{0.46\linewidth}
\begin{alignat*}{3}
    \max &\qquad&&{α} && \\
	\text{s.t.} &&& 0≤{α}≤\min_{{p}_i\in{P}}\braces*{\frac{{T}_i}{{e}_i}}&& \\
	&&& ∀ {p}_i\in{P}: && \sum_{{c}_m\in{C}}{a}_{i,m} = 1 \\
	&&&&& 0≤{t}_i≤{T}_i-{e}_i\\
    &&&&& ∀{c}_m∈\braces{{C}\smallsetminus{D}_i}:\ {a}_{i,m}=0\\
    &&& ∀ {c}_m\in{C}: && \sum_{{p}_i\in{P}}{a}_{i,m}{s}_i ≤  {S}_m \\
	&&& ∀ {p}_i,{p}_j\in{P}^2: && {f}_i ≠ {f}_j,\ ∀{c}_m\in{C}:\ {a}_{i,m}≤ {a}_{j,m} \\
	&&&&& {f}_i = {f}_j,\ ∀{c}_m\in{C}:\ {a}_{i,m}={a}_{j,m} \\
	&&&∀ {p}_i,{p}_j\in{P},j>i: && ∀{c}_m\in{C}:\ {t}_j-{t}_i-{q}_{i,j}{g}_{i,j}≥{α}{e}_i-\\
	&&&&&-Z\left(2-{a}_{i,m}-{a}_{j,m}\right) \\ 
\end{alignat*}
\end{minipage}
\hfill
\begin{minipage}{0.46\linewidth}
\begin{alignat*}{3}
	&&&&& ∀{c}_m\in{C}:\ {t}_j-{t}_i-{q}_{i,j}{g}_{i,j}≤-{g}_{i,j}{α}{e}_j-\\
	&&&&&-Z\left(2-{a}_{i,m}-{a}_{j,m}\right) \\
	&&&&& \frac{{e}_i-{T}_i}{{g}_{i,j}} ≤ {q}_{i,j} ≤ \frac{{T}_j-{e}_i}{{g}_{i,j}} \\
	&&&&& 0≤{t}_j-{t}_i-{q}_{i,j}{g}_{i,j}≤{g}_{i,j} \\
	&&&∀ {p}_i,{p}_j\in{P},{E_{ij}^{max}}<∞: && {t}_j-{t}_i-{q}_{i,j}{g}_{i,j}+{e}_j+{x}_{i,j}{T}_j≤{E_{ij}^{max}} \\
	&&&&& {t}_j-{t}_i-{q}_{i,j}{g}_{i,j}+{e}_i-\sum_{{c}_m,{c}_n\in{C}}\left({y}_{i,j,m,n}{τ}_{m,n}\right) + \\
	&&&&& + {x}_{i,j}Z ≥ 0 \\
	&&&&& ∀{c}_m,{c}_n\in{C}:\ {y}_{i,j,m,n}≥ {a}_{i,m}+{a}_{j,n}-1 \\
	&&&&& ∀{c}_m,{c}_n\in{C}:\ {y}_{i,j,m,n}≤ {a}_{i,m} \\
	&&&&& ∀{c}_m,{c}_n\in{C}:\ {y}_{i,j,m,n}≤ {a}_{j,n} \\
	&&& {a}_{i,m}\in\braces*{0,1},\ {t}_i\in\mathbb{Z},\ {q}_{i,j}\in\mathbb{Z}&\;& \\
	&&& {x}_{i,j}\in\braces*{0,1},\ {y}_{i,j,m,n}\in\braces*{0,1}&&
\end{alignat*}
\end{minipage}
}}
    \note<1->{
        \begin{itemize}
            \item formalismo
            \item garante optimalidade
            \item decisão é NP-completa, problema de optimalidade é NP-dificil (literatura)
            \item P polinomial, NP verificacao polinomial, NP-hard NP reduzidos a em tempo polinomial, NP-complete intersecao de NP e NP-hard
        \end{itemize}
        \centering
        \resizebox{!}{4.5cm}{\includegraphics{graphics/np_squeme.png}}
    }
\end{frame}

\subsection{Heuristic optimization}

\begin{frame}
    \frametitle{Best Response Algorithm \cite{al2012strictly}}

    \begin{itemize}
        \item<1-> Based on the Game Theory algorithm of the same name
        \item<2-> Partitions sequentially move to the best available offset
        \item<3-> When no partition can improve its $α$ value, equilibrium is reached
        \item<4-> Optimality dependant on the starting point
        \item<5-> Time complexity: $\bigO(N·T^{N-1})$
    \end{itemize}

    \vfill

    {\centering\resizebox{0.95\linewidth}{!}{\begin{tikzpicture}[scale=1.0]
	%\tikzstyle my grid=[style=help lines, color=blue!50, step=1, very thin, dashed]
	%\draw[style=my grid] (-1, -0.5) grid (26,6);	

	\pgfmathsetmacro\h{0.5};

    \foreach \x in {0, 24} {
        \draw[->, swred] (\x, 0.5) -- (\x, 5) node[above, black] {\scriptsize MTF};
    }
    \foreach \x in {1,...,23} {
        \draw[dashed, gray] (\x, 0.5) -- (\x, 5);
    }
	
	\pgfmathsetmacro\y{4};
    \pgfmathsetmacro\w{3};
	\foreach \x in {0, 12} {
		\pgfmathsetmacro\r{\x+\w}
		\pgfmathsetmacro\t{\y+\h};
		\draw[gray, fill=colorA] (\x, \y) rectangle node[black] {$p_1$} (\r, \t);
	}
	\pgfmathsetmacro\y{4};
    \pgfmathsetmacro\w{2};
	\foreach \x in {10} {
		\pgfmathsetmacro\r{\x+\w}
		\pgfmathsetmacro\t{\y+\h};
		\draw[gray, fill=colorB] (\x, \y) rectangle node[black] {$p_2$} (\r, \t);
	}
	\pgfmathsetmacro\y{4};
    \pgfmathsetmacro\w{1};
	\foreach \x in {14} {
		\pgfmathsetmacro\r{\x+\w}
		\pgfmathsetmacro\t{\y+\h};
		\draw[gray, pattern=crosshatch, pattern color=colorC] (\x, \y) rectangle node[black] {$p_3$} (\r, \t);
	}
    \pgfmathsetmacro\l{\y+0.5*\h}
    \draw (0, \l) node[left] {a)};



	\pgfmathsetmacro\y{3};
    \pgfmathsetmacro\w{3};
	\foreach \x in {0, 12} {
		\pgfmathsetmacro\r{\x+\w}
		\pgfmathsetmacro\t{\y+\h};
		\draw[gray, fill=colorA] (\x, \y) rectangle node[black] {$p_1$} (\r, \t);
	}
	\pgfmathsetmacro\y{3};
    \pgfmathsetmacro\w{2};
	\foreach \x in {7} {
		\pgfmathsetmacro\r{\x+\w}
		\pgfmathsetmacro\t{\y+\h};
		\draw[gray, fill=colorB] (\x, \y) rectangle node[black] {$p_2$} (\r, \t);
	}
	\pgfmathsetmacro\y{3};
    \pgfmathsetmacro\w{1};
	\foreach \x in {14} {
		\pgfmathsetmacro\r{\x+\w}
		\pgfmathsetmacro\t{\y+\h};
		\draw[gray, pattern=crosshatch, pattern color=colorC] (\x, \y) rectangle node[black] {$p_3$} (\r, \t);
	}
    \pgfmathsetmacro\l{\y+0.5*\h}
    \draw (0, \l) node[left] {b)};


	\pgfmathsetmacro\y{2};
    \pgfmathsetmacro\w{3};
	\foreach \x in {0, 12} {
		\pgfmathsetmacro\r{\x+\w}
		\pgfmathsetmacro\t{\y+\h};
		\draw[gray, fill=colorA] (\x, \y) rectangle node[black] {$p_1$} (\r, \t);
	}
	\pgfmathsetmacro\y{2};
    \pgfmathsetmacro\w{2};
	\foreach \x in {7} {
		\pgfmathsetmacro\r{\x+\w}
		\pgfmathsetmacro\t{\y+\h};
		\draw[gray, fill=colorB] (\x, \y) rectangle node[black] {$p_2$} (\r, \t);
	}
	\pgfmathsetmacro\y{2};
    \pgfmathsetmacro\w{1};
	\foreach \x in {21} {
		\pgfmathsetmacro\r{\x+\w}
		\pgfmathsetmacro\t{\y+\h};
		\draw[gray, fill=colorC] (\x, \y) rectangle node[black] {$p_3$} (\r, \t);
	}
    \pgfmathsetmacro\l{\y+0.5*\h}
    \draw (0, \l) node[left] {c)};


	\pgfmathsetmacro\y{1};
    \pgfmathsetmacro\w{3};
	\foreach \x in {0, 12} {
		\pgfmathsetmacro\r{\x+\w}
		\pgfmathsetmacro\t{\y+\h};
		\draw[gray, fill=colorA] (\x, \y) rectangle node[black] {$p_1$} (\r, \t);
	}
	\pgfmathsetmacro\y{1};
    \pgfmathsetmacro\w{2};
	\foreach \x in {7} {
		\pgfmathsetmacro\r{\x+\w}
		\pgfmathsetmacro\t{\y+\h};
		\draw[gray, fill=colorB] (\x, \y) rectangle node[black] {$p_2$} (\r, \t);
	}
	\pgfmathsetmacro\y{1};
    \pgfmathsetmacro\w{1};
	\foreach \x in {21} {
		\pgfmathsetmacro\r{\x+\w}
		\pgfmathsetmacro\t{\y+\h};
		\draw[gray, fill=colorC] (\x, \y) rectangle node[black] {$p_3$} (\r, \t);
	}
    \pgfmathsetmacro\l{\y+0.5*\h}
    \draw (0, \l) node[left] {d)};

    %%% utilities

	\pgfmathsetmacro\y{3};
    \pgfmathsetmacro\w{4};
    \pgfmathsetmacro\x{3}
    \pgfmathsetmacro\r{\x+\w}
    \pgfmathsetmacro\t{\y+\h};
    \draw[gray, dashed, fill=colorA, fill opacity=0.3] (\x, \y) rectangle (\r, \t);
    \draw[<->] (3, 3.6) -- node[above] {$(1-α_2)e_1$} (7, 3.6);

	\pgfmathsetmacro\y{3};
    \pgfmathsetmacro\w{2.666};
    \pgfmathsetmacro\x{9}
    \pgfmathsetmacro\r{\x+\w}
    \pgfmathsetmacro\t{\y+\h};
    \draw[gray, dashed, fill=colorB, fill opacity=0.3] (\x, \y) rectangle (\r, \t);
    \draw[<->] (9, 2.9) -- node[below] {$(1-α_2)e_2$} (11.6666, 2.9);

	\pgfmathsetmacro\y{2};
    \pgfmathsetmacro\w{6};
    \pgfmathsetmacro\x{15}
    \pgfmathsetmacro\r{\x+\w}
    \pgfmathsetmacro\t{\y+\h};
    \draw[gray, dashed, fill=colorA, fill opacity=0.3] (\x, \y) rectangle (\r, \t);
    \draw[<->] (15, 2.6) -- node[above] {$(1-α_3)e_1$} (21, 2.6);

	\pgfmathsetmacro\y{2};
    \pgfmathsetmacro\w{2};
    \pgfmathsetmacro\x{22}
    \pgfmathsetmacro\r{\x+\w}
    \pgfmathsetmacro\t{\y+\h};
    \draw[gray, dashed, fill=colorC, fill opacity=0.3] (\x, \y) rectangle (\r, \t);
    \draw[<->] (21, 2.6) -- node[above] {$α_3 e_3$} (24, 2.6);

\end{tikzpicture}
}}
    \note<1->{
        \begin{itemize}
            \item Achar uma solução melhor, (possivelmente) sub-óptima, em tempo razoável
        \end{itemize}
    }
\end{frame}

\subsection{Global Optimization}

\begin{frame}
    \frametitle{Global Optimization}

    Explore different distributions of partitions among modules
    \begin{itemize}
        \item Optimize each module's schedule using the previous algorithms
        \item Change the distribution and repeat
    \end{itemize}
    \note<1->{
        \begin{itemize}
            \item Métodos anteriores optimizam um módulo
        \end{itemize}
    }
\end{frame}

\begin{frame}
    \frametitle{Stochastic Optimization Algorithms}

    Operators:
    \begin{itemize}
        \item Move partitions between two modules
        \item Swap partitions between two modules
        \item Shuffle partition offsets
        \item Add/remove execution windows
    \end{itemize}

    \vspace{15pt}

    Meta-heuristic algorithms:
    \begin{itemize}
        \item Simulated Annealing
        \item Tabu Search
        \item Genetic Algorithm
    \end{itemize}
    % mention Stopping criterion...
    \note<1->{
        \begin{itemize}
            \item Algoritmos bastante estudados na literatura
            \item Parar quando temos uma solução tão boa como desejada, ou ao fim de um tempo limite
        \end{itemize}
    }
\end{frame}

% Not sure if I will add this slide
%\begin{frame}
%    \frametitle{Tabu Search}
%
%    Evaluate neighborhood of the current solution
%    Accept the best solution which was not recently visited
%
%\end{frame}

\section{Results}

\begin{frame}
    \frametitle{Partition Scheduling Model}
    \tableofcontents[currentsection,hideallsubsections]
\end{frame}

\subsection{Algorithmic Performance}

\begin{frame}
    \frametitle{Best Response Algorithm}

    \resizebox{!}{6cm}{\begin{tikzpicture}
\begin{axis}[
        ylabel={$t$ [\si{\milli\second}]},
        ylabel style={rotate=-90},
        xlabel={$N_p$},
        ymajorgrids,
    ]

    \addplot[
        color=colorA,
        mark=*,
    ] table [col sep = comma, x index = 0, y index = 1] {best-response-by-partition.csv};

\end{axis}
\end{tikzpicture}
}
    \note<1->{
        \begin{itemize}
            \item Em relação aos algoritmos relacionados com a Melhor Resposta
        \end{itemize}
    }
\end{frame}

\subsection{Scheduling Tool}

\begin{frame}
    \frametitle{Test Cases}

    \centering
    \begin{tabular}{c c c c r}
        \toprule
        Designation & Modules & Partitions & Chains & $α_{best}$ \\
        \midrule
        $2M6P$     & \num{2}  & \num{6}   & \num{0}  & \num{5.5}*   \\
        $4M10P$    & \num{4}  & \num{10}  & \num{3}  & \num{6.403}  \\ 
        $4M20P$    & \num{4}  & \num{20}  & \num{8}  & \num{2.875}  \\ 
        $8M40P$    & \num{8}  & \num{40}  & \num{15} & \num{2.984}  \\
        $20M100P$  & \num{20} & \num{100} & \num{40} & \num{2.325}  \\
        \midrule
        $3M15P\mathdash S$ & \num{3} & \num{15} & \num{3} & \num{1.26} \\
        \bottomrule
    \end{tabular}

    \vspace{25pt}

    $T∈\braces*{100,200,500,1000}$

    \note<1->{
        \begin{itemize}
            \item problema 20M100P tem dimensão semelhante aos enconctrados na industria
            \item periodos não harmónicos
            \item Reforçar, Divisao em varias janelas tratada separadamente
        \end{itemize}
    }
\end{frame}

\begin{frame}
    \frametitle{Feasible Solution}
    \centering
    \begin{tabular}{c c c}
        \toprule
        Instance           & $t_{MILP}$          & $t_{heuristic}$ (median) \\
        \midrule
        $2M6P$             & \SI{1.00}{\second}  & \SI{0.593}{\second} \\
        $4M10P$            & \SI{1.34}{\second}  & \SI{0.595}{\second} \\ 
        $4M20P$            & \SI{6.29}{\second}  & \SI{0.830}{\second} \\ 
        $8M40P$            & \SI{310.7}{\second} & \SI{1.155}{\second} \\
        $20M100P$          & $>\SI{24}{\hour}$   & \SI{23.32}{\second} \\
        \midrule
        $3M15P\mathdash S$ & NA                  & \SI{74.53}{\second} \\
        \bottomrule
    \end{tabular}
\end{frame}

\begin{frame}
    \frametitle{Improved solution}

    MILP solver does not converge in under \SI{24}{\hour}, except for $2M6P$.

    \vfill

    \begin{minipage}{0.48\linewidth}
        \centering
        $2M6P$
        \vspace{10pt}

        \resizebox{0.9\linewidth}{!}{\begin{tikzpicture}
\begin{axis}[
	boxplot/draw direction=x,
	ytick={1, 2, 3},
	yticklabels={SA, Tabu, Genetic},
	xlabel={Time to $\alpha_{best}$ [\si{\second}]},
    width=10.6cm,
    height=4.5cm,
    scale only axis=true,
	]
	
\addplot+ [
colorA,
thick,
mark options={fill=colorA},
boxplot prepared={
lower whisker=0.5600,
lower quartile=0.5800,
median=0.5950,
upper quartile=0.6225,
upper whisker=0.6800,
},] coordinates {};

\addplot+ [
colorC,
thick,
mark options={fill=colorC},
boxplot prepared={
lower whisker=0.5900,
lower quartile=0.6000,
median=0.6600,
upper quartile=0.7375,
upper whisker=0.8200,
},] coordinates {(0, 0.99)};

\addplot+ [
colorD,
thick,
mark options={fill=colorD},
boxplot prepared={
lower whisker=0.5800,
lower quartile=0.6300,
median=0.6700,
upper quartile=0.7625,
upper whisker=0.8200,
},] coordinates {};

\end{axis}
\end{tikzpicture}
}
    \end{minipage}
    \hfill
    \begin{minipage}{0.48\linewidth}
        \centering
        $2M10P$
        \vspace{10pt}

        \resizebox{0.9\linewidth}{!}{\begin{tikzpicture}
\begin{axis}[
	boxplot/draw direction=x,
	ytick={1, 2, 3},
	yticklabels={SA, Tabu, Genetic},
	xlabel={Time to $\alpha_{best}$ [\si{\second}]},
	]
	
\addplot+ [
colorA,
thick,
mark options={fill=colorA},
boxplot prepared={
lower whisker=1.72,
lower quartile=7.6125,
median=15.969999999999999,
upper quartile=24.0325,
upper whisker=35.57,
},] coordinates {};

\addplot+ [
colorC,
thick,
mark options={fill=colorC},
boxplot prepared={
lower whisker=2.2199999999999998,
lower quartile=6.66,
median=17.83,
upper quartile=47.855000000000004,
upper whisker=69.55,
},] coordinates {};

\addplot+ [
colorD,
thick,
mark options={fill=colorD},
boxplot prepared={
lower whisker=7.08,
lower quartile=16.97,
median=24.08,
upper quartile=43.22,
upper whisker=71.63,
},] coordinates {};

\end{axis}
\end{tikzpicture}
}
    \end{minipage}
    \note<1->{
        \begin{itemize}
            \item casos de estudo corridos várias vezes, dado serem estocásticos
            \item Por conveniência, pára-se numa solução proxima da melhor
        \end{itemize}
    }
\end{frame}

\begin{frame}
    \frametitle{Improved solution}

    \begin{minipage}{0.48\linewidth}
        \centering
        $4M20P$
        \vspace{10pt}

        \resizebox{0.9\linewidth}{!}{\begin{tikzpicture}
\begin{axis}[
	boxplot/draw direction=x,
	ytick={1, 2, 3},
	yticklabels={SA, Tabu, Genetic},
	xlabel={Time to $0.9\alpha_{best}$ [\si{\minute}]},
	]
	
\addplot+ [
colorA,
thick,
mark options={fill=colorA},
boxplot prepared={
lower whisker=0.2063,
lower quartile=0.9390,
median=2.5591,
upper quartile=4.0567,
upper whisker=8.6253,
},] coordinates {(0, 10.335833333333333)};

\addplot+ [
colorC,
thick,
mark options={fill=colorC},
boxplot prepared={
lower whisker=0.9623,
lower quartile=1.5548,
median=2.1759,
upper quartile=7.0902,
upper whisker=11.2707,
},] coordinates {(0, 19.3535)};

\addplot+ [
colorD,
thick,
mark options={fill=colorD},
boxplot prepared={
lower whisker=2.1847,
lower quartile=5.1395,
median=6.9888,
upper quartile=11.2683,
upper whisker=12.8303,
},] coordinates {};

\end{axis}
\end{tikzpicture}
}
    \end{minipage}
    \hfill
    \begin{minipage}{0.48\linewidth}
        \centering
        $8M40P$
        \vspace{10pt}

        \resizebox{0.9\linewidth}{!}{\begin{tikzpicture}
\begin{axis}[
	boxplot/draw direction=x,
	ytick={1, 2, 3},
	yticklabels={SA, Tabu, Genetic},
	xlabel={Time to $0.9\alpha_{best}$ [\si{\minute}]},
	]
	
\addplot+ [
colorA,
thick,
mark options={fill=colorA},
boxplot prepared={
lower whisker=5.6307,
lower quartile=11.4145,
median=12.8042,
upper quartile=16.1277,
upper whisker=22.9553,
},] coordinates {(0, 28.1835)};

\addplot+ [
colorC,
thick,
mark options={fill=colorC},
boxplot prepared={
lower whisker=3.8188,
lower quartile=7.5687,
median=11.1314,
upper quartile=13.3697,
upper whisker=15.9153,
},] coordinates {(0, 22.764166666666664) (0, 85.81333333333333)};

\addplot+ [
colorD,
thick,
mark options={fill=colorD},
boxplot prepared={
lower whisker=26.2935,
lower quartile=32.6950,
median=35.8912,
upper quartile=38.8532,
upper whisker=45.0547,
},] coordinates {};

\end{axis}
\end{tikzpicture}
}
    \end{minipage}

\end{frame}

\begin{frame}
    \frametitle{Improved solution}

    \begin{minipage}{0.48\linewidth}
        \centering
        $20M100P$
        \vspace{10pt}

        \resizebox{0.9\linewidth}{!}{\begin{tikzpicture}
\begin{axis}[
	boxplot/draw direction=x,
	ytick={1, 2},
	yticklabels={SA, Tabu},
	xlabel={Time to $0.9\alpha_{best}$ [\si{\minute}]},
	]
	
\addplot+ [
colorA,
thick,
mark options={fill=colorA},
boxplot prepared={
lower whisker=11.8787,
lower quartile=40.1135,
median=60.3086,
upper quartile=88.5171,
upper whisker=120.4787,
},] coordinates {};

\addplot+ [
colorC,
thick,
mark options={fill=colorC},
boxplot prepared={
lower whisker=17.5713,
lower quartile=30.7858,
median=42.4836,
upper quartile=56.1466,
upper whisker=73.9508,
},] coordinates {(0, 104.76616666666668) (0, 110.62766666666667)};

\end{axis}
\end{tikzpicture}
}
    \end{minipage}
    \hfill
    \begin{minipage}{0.48\linewidth}
        \centering
        $3M15P\mathdash S$
        \vspace{10pt}

        \resizebox{0.9\linewidth}{!}{\begin{tikzpicture}
\begin{axis}[
	boxplot/draw direction=x,
	ytick={1, 2, 3},
	yticklabels={SA, Tabu, Genetic},
	xlabel={Time to $\alpha_{best}$ [\si{\minute}]},
    xtick={0, 10, 20, 30, 40, 50, 60},
    xticklabels={$0$, $10$, $20$, $30$, $40$, $50$,},
	]

\addplot+ [
colorA,
thick,
mark options={fill=colorA},
boxplot prepared={
lower whisker=0.9343,
lower quartile=16.3835,
median=31.6933,
upper quartile=38.6162,
upper whisker=47.4082,
},] coordinates {};

\addplot+ [
colorC,
thick,
mark options={fill=colorC},
boxplot prepared={
lower whisker=0.5437,
lower quartile=8.3616,
median=12.6054,
upper quartile=18.9163,
upper whisker=22.2335,
},] coordinates {(0, 40.369166666666665) (0, 43.35783333333333)};

\addplot+ [
colorD,
thick,
boxplot prepared={
lower whisker=1.5002,
lower quartile=7.3434,
median=13.4660,
upper quartile=21.5471,
upper whisker=36.0757,
},] coordinates {(0, 45.88) (0, 54.65133333333333) (0, 55.0815)};

\end{axis}
\end{tikzpicture}
}
    \end{minipage}
    \note<1->{
        \begin{itemize}
            \item Uma hora parece muito, mas no tempo que estes sistemas demoram a serem integrados, não é nada
            \item Relembrar que temos solução possível em segundos
        \end{itemize}
    }
\end{frame}


\section{Conclusions}

\begin{frame}
    \frametitle{Achievements}

    \begin{itemize}
        \item<1-> A scheduling tool and framework that supports a variety of constraints
        \vspace{10pt}
        \item<2-> Capable of producing flexible solutions in moderate amounts of time
        \vspace{10pt}
        \item<3-> A mathematical model that yields optimal solutions when time is not an issue
    \end{itemize}
    \note<1->{
        \begin{itemize}
            \item mencionar GMV
        \end{itemize}
    }
\end{frame}

\begin{frame}
    \frametitle{Recommendations for Future Work}

    \begin{itemize}
        \item<1-> Integration with other (D)IMA configuration or V\&V frameworks
        \vspace{30pt}
        \item<2-> Partition scheduling in multicore (D)IMA systems
    \end{itemize}
    \note<1->{
        \begin{itemize}
            \item Configuracao de redes, decisões acerca do numero e tipo de processadores
            \item Multicore traz nuances que nao sao abordadas aqui
        \end{itemize}
    }
\end{frame}

\nocite{samolej2011arinc, pira2016line}

\begin{frame}
    \frametitle{Selected Bibliography}
    \printbibliography
\end{frame}

\end{document}
